\section{非退化性与辛位移能}
\label{sec:nondegeneracy-displacement-energy}

本节先说明霍费尔 (H90) 所采用的伪范数非退化性证明思路。

下面是霍费尔在 $\mathbb C^n$ 上证明非退化性时使用的关键概念。

\begin{definition}[位移能]\label{def:displacement-energy}
\stfhIndex{I102}{wei^yin^eng}{位移能（displacement energy）}
设 $A\subset M$ 是 $M$ 的任意子集。定义
\begin{equation}
 e(A)=\inf_H\{\|H\|\mid
   \overline A\cap\phi_H^1(A)=\varnothing\},
 \label{eq:displacement-energy}
\end{equation}
并称 $e(A)$ 为 $A$ 的\emph{位移能}。
\end{definition}

显然，位移能在 $\Symp(M,\omega)$ 的作用下不变。

假设 $\phi:M\to M$ 不是恒等映射。则存在一点 $x\in M$，使得
$\phi(x)\ne x$。由 $\phi$ 的连续性，确实可以找到 $x$ 的一个开邻域
$U$，使得
\[
 \phi(U)\cap\overline U=\varnothing.
\]
由达布定理，可选取 $\overline U$，使它是某个辛嵌入
\[
 g:B^{2n}(r)\subset\mathbb C^n\longrightarrow(M,\omega);
 \qquad g(0)=x
\]
\stfhIndex{I142}{ge^luo^mo^fu^mian^ji}{格罗莫夫面积（Gromov area）}
的像。一般地，把任意这类辛嵌入的像称为格罗莫夫面积
$\lambda=\pi r^2$ 的\emph{辛球}。以 $B(\lambda)$ 表示任意这类格罗莫夫面积为 $\lambda$ 的\emph{辛球}。
\stfhIndex{I271}{xin^qiu}{辛球（symplectic ball）}

注意，所有满足 $\|\phi\|=0$ 的“零”哈密顿微分同胚 $\phi$ 组成一个正规子群。因此，巴尼亚加定理蕴含零哈密顿微分同胚的集合或者是
$\{\operatorname{id}\}$，或者是整个群 $\Ham(M,\omega)$。所以只需证明：对
$\lambda>0$ 的辛球 $B(\lambda)\subset M$，有 $e(B(\lambda))>0$；因为根据位移能的定义，这随即蕴含存在一个哈密顿微分同胚
$\operatorname{id}\ne\phi$，满足
$\|\phi\|\geq e(B(\lambda))>0$。

为了表明 $e(B(\lambda))>0$ 的正性证明并非平凡，我们陈述一个命题；它说明把 $M$ 中的一点移到任意另一点只需“零能量”。

\begin{proposition}\label{prop:point-zero-displacement-energy}
设 $x_0,x_1$ 是 $M$ 中任意两点。则对任意给定的 $\delta>0$，存在
$H$，使得 $\phi_H^1(x_0)=x_1$ 且 $\|H\|\leq\delta$。换言之，一点的位移能为零。
\end{proposition}

\begin{proof}
先从 $M=\mathbb R^{2n}$ 且带典范辛形式的情形开始。注意辛群
$Sp(2n,\mathbb R)$ 在 $\mathbb R^{2n}$ 上传递作用，因而可以假设
$x_0=0$ 为原点，而对某个 $s\in\mathbb R$，
$x_1=(s,0,\ldots,0)$。

其次，考虑自治哈密顿量 $h=sp_1$。则
$X_{sp_1}=s\,\partial/\partial q^1$，其流由
\[
 \phi_{sp_1}^t(x)=x+t(s,\ldots,0),
 \qquad x\in\mathbb R^{2n}
\]
给出。特别地，$\phi_h^1(x_0)=x_1$。注意，对所有 $t$，都有
$p_1(\phi_h^t(x_0))=0$。

最后，考虑线段 $I:=\overline{x_0x_1}$，并以 $N_\delta(I)$ 表示其
$\delta$-邻域。选取截断函数 $\rho:\mathbb R^{2n}\to\mathbb R$，满足
\[
 \rho(x)=
 \begin{cases}
  1,&x\in N_{\delta/3}(I),\\
  0,&x\notin N_{\delta/2}(I),
 \end{cases}
 \qquad 0\leq\rho(x)\leq1.
\]
然后选取哈密顿量 $H=\rho\cdot h$。因为在
$N_{\delta/2}(I)\supset I$ 上有 $H\equiv h$，所以
$\phi_H^1(x_0)=\phi_h^1(0)=x_1$。此外，由 $H$ 的构造可得
\[
 \max H\leq\frac\delta2,
 \qquad
 \min H\geq-\frac\delta2,
\]
从而 $\|H\|\leq\delta$。由于 $\delta>0$ 可任意小，这就完成了
$\mathbb R^{2n}$ 情形的证明。

现在考虑一般的 $(M,\omega)$。任选嵌入曲线
$\alpha:[0,1]\to M$，满足 $\alpha(0)=x_0$、$\alpha(1)=x_1$。选取网格为
$1/N$ 的有限分割
\[
 0=t_0<t_1<\cdots<t_N=1.
\]
通过把 $N$ 取得充分大，可以选取 $\alpha([0,1])$ 的一个由一致大小的有限多个辛球
\[
 g_i:B^{2n}(r)\to M,
 \qquad i=0,\ldots,N-1
\]
组成的覆盖，使得每段 $\alpha([t_i,t_{i+1}])$ 都包含在该覆盖中的一个球内。只需证明，对每个 $i=0,\ldots,N-1$，把点
$\alpha(t_i)$ 移到 $\alpha(t_{i+1})$ 只需“零能量”。然而，这来自证明第一部分中 $\mathbb R^{2n}$ 上的结果：把第一部分应用于
$B^{2n}(r)\subset\mathbb R^{2n}$，再把相应哈密顿量乘以适当截断函数后推前。这就完成证明。
\end{proof}

\begin{exercise}\label{ex:point-displacement-proof-details}
补全上述证明论证最后一部分的细节。更准确地说，证明可以选取
$\alpha([0,1])$ 的一个由有限多个一致大小的辛球组成的覆盖，并详细说明如何通过 $g_i$ 把 $B^{2n}(r)$ 上的结果转移到 $(M,\omega)$。
\end{exercise}

另一方面，下面这个非平凡定理说明，任意 $\lambda>0$ 的辛球
$B(\lambda)$ 的位移能为正。

\begin{theorem}\label{thm:energy-capacity-basic}
设 $B(\lambda)$ 是任意辛球，$\lambda>0$。
\begin{enumerate}[label=(\arabic*)]
 \item \textup{\bfseries（霍费尔 (H90)）}设 $B^{2n}(r)$ 是 $\mathbb R^{2n}$ 中的标准球。则
 $e(B^{2n}(r))=\pi r^2$。
 \item \textup{\bfseries（拉隆德--麦克达夫 (LM95a)）}设 $B(\lambda)\subset M$ 是任意格罗莫夫面积为 $\lambda>0$ 的辛球，则
 $e(B(\lambda))\geq\lambda/2$。
\end{enumerate}
\end{theorem}

这个定理有时称为\emph{能量--容量不等式}。厄舍 (Ush10a) 证明了下面这个能量--容量不等式的最优版本。我们将在第 22.3 节给出其证明。

\begin{theorem}[最优能量--容量不等式，(Ush10a)]
\stfhIndex{I221}{zui^you^neng^liang^rong^liang^bu^deng^shi}{最优能量--容量不等式（optimal energy–capacity inequality）}
\stfhIndex{I297}{e^she}{厄舍（Usher）}
\label{thm:optimal-energy-capacity}
设 $(M,\omega)$ 是任意辛流形，且 $B(\lambda)\subset M$ 是一个辛球。则
\[
 e(B(\lambda))\geq\lambda.
\]
\end{theorem}

$\lambda>0$ 时 $e(B(\lambda))>0$ 的正性结论立即给出如下推论。

\begin{corollary}\label{cor:hofer-norm-nondegenerate}
$\Ham(M,\omega)$ 上的霍费尔范数是非退化的；即
$\|\phi\|=0$ 当且仅当 $\phi=\operatorname{id}$。
\end{corollary}

需要提到，霍费尔距离
$d(\phi,\phi'):=\|\phi^{-1}\phi'\|$ 是内蕴的，因为它正是由
$C_0^\infty(M)$ 上定义的伴随不变范数
\[
 \|h\|=\operatorname{osc}(h)=\max h-\min h
\]
诱导的芬斯勒度量之距离。显然，$\Ham(M,\omega)$ 上的这个霍费尔范数是双不变的，并且关于
$\Ham(M,\omega)\subset\Diff(M)$ 的 $C^\infty$ 拓扑连续。在这一方面，布霍夫斯基与奥斯特罗弗结合定理 5.2.4 与定理 5.2.3，证明了如下唯一性结果。

\begin{theorem}[布霍夫斯基--奥斯特罗弗]\label{thm:hofer-finsler-uniqueness}
设 $(M,\omega)$ 是闭的。则 $\Ham(M,\omega)$ 上任何由
$C_0^\infty(M)$ 上一个在 $C^\infty$ 拓扑中连续的伪范数
$|\mathord\cdot|$ 得到的双不变芬斯勒伪度量，或者恒等于零，或者等价于霍费尔的芬斯勒度量。特别地，任何非退化双不变芬斯勒度量都等价于霍费尔度量。
\end{theorem}

结束本节之前，我们想重新审视本小节开头概述的非退化性证明方案。从证明容易看出，除了辛球 $B(\lambda)$ 的集合以外，只要 $M$ 的子集族
$C$ 满足以下两个性质，也完全可以使用它：
\begin{itemize}
 \item 对任意开子集 $U\subset M$，存在 $A\in C$，使得
 $A\subset U$；
 \item 对任意 $A\in C$，都能证明 $e(A)>0$。
\end{itemize}
\stfhIndex{I144}{ge^luo^mo^fu^de^8^zi^xing^ji^qiao}{格罗莫夫的 8 字形技巧（Gromov’s figure-8 trick）}
\stfhIndex{I228}{bo^er^jie^luo^wei^qi}{波尔捷罗维奇（Polterovich）}
波尔捷罗维奇 (Po93) 观察到，拉格朗日子流形的集合（甚至拉格朗日环面的集合）满足这两个性质。使用拉格朗日子流形非常适合格罗莫夫的伪全纯曲线方法。我们将给出 (Oh97c) 中的一种简单证明：把对弗洛尔扰动柯西--黎曼方程的研究作一个容易的应用，证明驯服辛流形中任意紧拉格朗日子流形的位移能为正。这个正性首先由波尔捷罗维奇 (Po93) 对\emph{有理}拉格朗日子流形证明，由切卡诺夫 (Che98) 对一般拉格朗日子流形证明。波尔捷罗维奇在证明中使用了格罗莫夫的 8 字形技巧 (Gr85)，而切卡诺夫在证明中使用了相当复杂的弗洛尔同调理论。
